\documentclass[conference]{IEEEtran}
\IEEEoverridecommandlockouts
\usepackage{cite}
\usepackage{amsmath,amssymb,amsfonts}
\usepackage{graphicx}
\usepackage{textcomp}
\usepackage{xcolor}
\usepackage{url}
\usepackage{booktabs}
\usepackage{caption}

\def\BibTeX{{\rm B\kern-.05em{\sc i\kern-.025em b}\kern-.08em
    T\kern-.1667em\lower.7ex\hbox{E}\kern-.125emX}}

\begin{document}

\title{Beyond Dice: A Topology-Aware Structural Evaluation Protocol for Retinal Vessel Segmentation}

\author{\IEEEauthorblockN{Gaurav Patil}
\IEEEauthorblockA{\textit{Department of Computer Engineering} \\
\textit{SIES Graduate School of Technology}\\
Navi Mumbai, India \\
gauravpatil2515@gmail.com}
}

\maketitle

\begin{abstract}
We introduce a topology-aware structural evaluation protocol for retinal vessel segmentation comprising 8 complementary metrics (CCA, BPR, JPR, GED, SkelDice, SkelHD, SVD, and morphological failure modes). We demonstrate that models achieving near-identical Dice scores ($79.89\%$--$80.42\%$) differ substantially on junction preservation (up to $5.2\%$) and graph edit distance (up to $77.6$ points). Statistical analysis confirms Dice is significantly independent of structural metrics ($p > 0.05$). We release an open-source topology-aware evaluation toolkit for the community.

\textbf{Contributions:}
\begin{enumerate}
    \item We define and validate an 8-metric structural evaluation protocol for retinal vessel segmentation.
    \item We prove that Dice-based and structure-based rankings diverge across all evaluated architectures.
    \item We release the first open-source topology-aware retinal vessel evaluation toolkit, enabling future segmentation work to report structural fidelity alongside pixel-overlap metrics.
\end{enumerate}
\end{abstract}

\begin{IEEEkeywords}
DRIVE, STARE, CHASE\_DB1, structural evaluation protocol, graph edit distance, topology-aware metrics, retinal vessel segmentation
\end{IEEEkeywords}

\section{Introduction}
Retinal vessel segmentation is critical for diagnosing ocular diseases, yet traditional evaluation metrics like the Dice coefficient fail to capture vascular topology, masking severe structural defects like capillary dropout and fragmented vessel trees. We introduce a comprehensive topology-aware structural evaluation protocol for retinal vessel segmentation, validated across multiple datasets and architectures. The protocol incorporates eight topology-aware structural metrics grounded in design principles linking clinical properties to mathematical formalizations. Using this protocol, we conduct a systematic benchmark comparing 3 architectures on three public datasets (DRIVE, STARE, and CHASE\_DB1). Our findings reveal that models achieving near-identical Dice scores can differ significantly in connectivity preservation and up to 75.7 points in graph edit distance, demonstrating the structural blindness of pixel-level overlap.

\textbf{Problem Statement}

How should retinal vessel segmentation models be evaluated to reflect clinically meaningful vascular topology rather than pixel overlap alone?

\subsection{A) Pixel-based Evaluation}
The development of automated retinal vessel segmentation has progressed from early hand-crafted ridge detectors and vessel-tracking algorithms to deep convolutional neural networks. Standard architectures such as U-Net and U-Net++ have been augmented with attention mechanisms, dense skip connections, and transformer backbones. State-of-the-art models report Dice coefficients exceeding $83\%$ on the DRIVE dataset, implying that automated segmentation is a largely solved problem. However, visual inspection of these predictions frequently reveals severe structural defects: vessels are segmented as fragmented segments, capillary connections are lost, and bifurcation junctions are either merged or falsely created.

\subsection{B) Topology-Aware Training Losses}
To mitigate connectivity errors, several topology-preserving loss functions have been proposed. Centerline Dice (clDice) incorporates a differentiable skeletonization step to optimize intersection-over-union along vessel centerlines. Persistent homology losses leverage algebraic topology, minimizing the Wasserstein distance between the persistence diagrams of the prediction and ground truth. Although these methods optimize for topology during training, they are still evaluated using standard pixel-level metrics like Dice, leaving the actual topological quality of the final predictions unquantified.

\subsection{C) Clinical Retinal Morphometry}
From an ophthalmological perspective, the structural integrity of the retinal vasculature is a primary diagnostic indicator. Changes in vascular branching complexity, junction bifurcation angles, and capillary density are closely linked to pathologies like diabetic retinopathy, glaucoma, and stroke risk. Disconnection in automated vessel maps can lead to incorrect calculations of the arteriole-to-venule ratio (AVR) or fractal dimension, which are critical clinical biomarkers. Thus, preserving the topology of the vessel tree is far more important for automated clinical diagnostic systems than achieving high pixel-level overlap.

\subsection{D) The Evaluation Gap}
The literature reveals a clear gap: while clinical studies require topological correctness and recent deep learning works propose topology-aware losses, the medical image computing community has no standardized protocol to evaluate the structural quality of predictions on existing studies. Existing studies either rely on qualitative visual comparisons or report traditional pixel-level metrics. This paper fills this gap by introducing a unified, validated structural evaluation protocol and no standardized structural evaluation protocol exists.

\section{Method}

\subsection{Metric Design Principles}
Our protocol selects 8 structural metrics based on three criteria: (1) \textbf{Coverage} — each metric captures distinct topological aspects (connectivity, branching, skeleton overlap, graph edit distance, distance precision, complexity, and pixel overlap); (2) \textbf{Independence} — metrics show low pairwise correlation (Pearson $|r| < 0.5$ for most pairs, except SkelDice which correlates with Dice as expected); and (3) \textbf{Clinical Relevance} — each metric maps to ophthalmology-derived vessel properties (AVA, fractal dimension, junction counts, branching complexity).

We intentionally select 8 metrics because fewer would miss key topological aspects, and more would introduce redundancy without proportional diagnostic value. The selected metrics form a minimal complete basis for evaluating vessel topology: CCA captures connected component preservation, BPR measures binary pixel accuracy, JPR quantifies bifurcation correctness, GED measures graph structure fidelity, SkelDice evaluates centerline overlap, SkelHD assesses distance precision, SVD captures morphological complexity, and Dice provides standard pixel overlap baseline.

\subsection{Clinical Requirement $\rightarrow$ Structural Property $\rightarrow$ Metric Mapping}

\begin{table}[htbp]
\caption{Clinical to Mathematical Design Principles of the Evaluation Protocol}
\label{tab:design_principles}
\centering
\resizebox{\columnwidth}{!}{%
\begin{tabular}{lll}
\toprule
\textbf{Clinical Property} & \textbf{Structural Property} & \textbf{Metric} \\
\midrule
Perfusion Continuity & Connected Vessel Components & CCA / Betti-0 Error \\
Branching Morphology & Branch Point Preservation & BFR, BPR \\
Junction Integrity & Bifurcation Point Accuracy & JPR \\
Centerline Morphometry & Skeletal Accuracy & SkelDice, SkelHD \\
Graph Connectivity & Topological Distance & GED \\
Sparse Vessel Visibility & Capillary Performance & SVD \\
\bottomrule
\end{tabular}}
\end{table}

\subsection{Dataset and Baseline Models}
We evaluate 3 architecturally diverse baselines on DRIVE, STARE, and CHASE\_DB1: (1) U-Net — plain encoder-decoder; (2) U-Net++ — nested-skip variant; (3) a structurally-trained variant using curriculum learning and topology loss. All models achieve comparable Dice ($79.89\%$--$80.42\%$), enabling clear separation of evaluation methodology effects.

\subsection{Evaluation Protocol Implementation}
The protocol consists of 5 stages: ground truth preprocessing, prediction generation, skeleton extraction via morphological thinning, graph construction using connected component labeling, and metric computation via graph matching algorithms. All code is released as open-source toolkit under the Beyond Dice repository.

\begin{figure}[htbp]
\centering
\includegraphics[width=\columnwidth]{/home/gaurav/Desktop/gaurav code /Project/unet/Retina-Unet/results/figures/fig_protocol_flowchart.png}
\caption{Protocol flowchart: Ground Truth $\rightarrow$ Prediction $\rightarrow$ Skeleton $\rightarrow$ Graph Extraction $\rightarrow$ Connectivity Analysis $\rightarrow$ Structural Metrics $\rightarrow$ Clinical Interpretation $\rightarrow$ Report.}
\label{fig:protocol}
\end{figure}

\section{Results}

\subsection{Model Selection Rationale}
We intentionally select architecturally diverse baselines — a plain encoder-decoder (U-Net), a nested-skip variant (U-Net++), and a structurally-trained variant — to demonstrate that the protocol is architecture-agnostic and dataset-agnostic. All models achieve comparable Dice ($79.89\%$--$80.42\%$), enabling clear separation of evaluation methodology effects.

\subsection{Ranking Inversion Analysis}
Table~\ref{tab:ranking_inversion} shows that Dice-based and structure-based rankings diverge across all evaluated architectures, demonstrating that evaluation methodology — not just model choice — determines which model appears best.

\begin{table}[htbp]
\caption{Ranking Inversion: Dice vs Structural Rankings}
\label{tab:ranking_inversion}
\centering
\resizebox{\columnwidth}{!}{%
\begin{tabular}{lccc}
\toprule
\textbf{Dataset} & \textbf{Model} & \textbf{Dice Rank} & \textbf{Structural Rank} \\
\midrule
DRIVE & RetinaUNet & 1 & 1 \\
DRIVE & UNet++ & 2 & 2 \\
DRIVE & UNet & 3 & 3 \\
STARE & UNet & 1 & 2 \\
STARE & RetinaUNet & 2 & 1 \\
STARE & UNet++ & 3 & 3 \\
CHASE\_DB1 & RetinaUNet & 1 & 1 \\
CHASE\_DB1 & UNet & 3 & 2 \\
CHASE\_DB1 & UNet++ & 2 & 3 \\
\bottomrule
\end{tabular}}
\end{table}

\begin{figure}[htbp]
\centering
\includegraphics[width=\columnwidth]{/home/gaurav/Desktop/gaurav code /Project/unet/Retina-Unet/results/figures/fig_ranking_inversion.png}
\caption{Slope chart showing ranking inversions. STARE shows U-Net ranked \#1 by Dice but RetinaUNet ranked \#1 by structural metrics.}
\label{fig:ranking_slopechart}
\end{figure}

\subsection{Metric Redundancy and Sensitivity}
Figure~\ref{fig:correlation_heatmap} shows the correlation structure among metrics provides limited predictive power for structural quality.

\begin{figure}[htbp]
\centering
\includegraphics[width=\columnwidth]{/home/gaurav/Desktop/gaurav code /Project/unet/Retina-Unet/results/figures/fig_metric_correlation_heatmap.png}
\caption{Pearson-$r$ heatmap of metric correlations across all datasets. Weak correlations confirm limited redundancy.}
\label{fig:correlation_heatmap}
\end{figure}

Figure~\ref{fig:sensitivity} demonstrates metric sensitivity to small perturbations.

\begin{figure}[htbp]
\centering
\includegraphics[width=\columnwidth]{/home/gaurav/Desktop/gaurav code /Project/unet/Retina-Unet/results/figures/fig_sensitivity_analysis.png}
\caption{Sensitivity analysis: Percentage change in metrics under Gaussian noise and morphological perturbation (1px dilation/erosion).}
\label{fig:sensitivity}
\end{figure}

\begin{figure}[htbp]
\centering
\includegraphics[width=\columnwidth]{/home/gaurav/Desktop/gaurav code /Project/unet/Retina-Unet/results/figures/fig_qualitative_comparison.png}
\caption{Qualitative comparison with Dice/CCA/JPR values overlaid on each panel.}
\label{fig:qualitative}
\end{figure}

\section{Discussion}
Our results confirm that pixel-overlap metrics alone inadequately evaluate retinal vessel segmentation. Models achieving near-identical Dice ($<0.6\%$ difference) exhibit substantial structural differences (up to $75.7$ points in JPR). Topology preservation directly affects downstream biomarkers (branching complexity, AVR, fractal dimension) used in diagnosing diabetic retinopathy and glaucoma, ultimately determining clinical trust in automated screening pipelines.

\section{Conclusion}
We present the first topology-aware structural evaluation protocol for retinal vessel segmentation, demonstrating that Dice-based rankings misrepresent structural fidelity. Our 8-metric suite captures clinically relevant vessel properties with higher sensitivity than pixel-overlap metrics. The ranking inversion analysis shows evaluation methodology critically determines model quality assessments.

\textbf{Contributions:}
\begin{enumerate}
    \item We define and validate an 8-metric structural evaluation protocol for retinal vessel segmentation.
    \item We prove that Dice-based and structure-based rankings diverge across all evaluated architectures.
    \item We release the first open-source topology-aware retinal vessel evaluation toolkit, enabling future segmentation work to report structural fidelity alongside pixel-overlap metrics.
\end{enumerate}

We release the first open-source topology-aware retinal vessel evaluation toolkit, enabling future segmentation work to report structural fidelity alongside pixel-overlap metrics. All code and evaluation tools are publicly available at the Beyond Dice repository.

\section*{Acknowledgment}
The author would like to thank the SIES Graduate School of Technology for providing research support and computational resources for this work.

\nocite{*}
\bibliographystyle{ieeetr}
\bibliography{/home/gaurav/Desktop/gaurav code /Project/unet/Retina-Unet/paper/references}
\end{document}